Requisitos gerais para a divulgaação de Informações financeiras Relacionadas a Sustentabilidade
O que significa ``razoavelmente afetar as perspectivas da entidade''?
Relações como outras entidades em rede de relacionamento.

Preocupação com a efetividade(qualidade das informações) que apresentam o impacto das oportunidades e riscos da sustentabilidade no : fluxo de caixa; acesso a financiamento e custo de capital no curto médio e longo prazo.
Não se pode repostar fatos sobre outra entidade. Só a propria entidade tem competência para elaborar os relatórios com as informações financeiras Relacionadas a oportunidades e riscos a Sustentabilidade.
Conteúdo (tudos esse itens se referem a oportunidades e riscos a sustentabilidade: Governança (processos, controles e procedientos) ; estrtégia de gerenciamento; processo específico de gestão de riscos; metas e metricas.

Como descrever a governança:
     Que individuies ou orgão tratan do assunto: descrição de funçoes ; mandatos e políticas
     Quais as habilidades e competências são exigitas para os indivíduos ou orgão e o procedimento de verificação se os reguisitos são cumpridos.
     Como e com que frequencia os orgãos ou individuos são informados dos riscos e oportunidades.
     Como os riscos e oportunidades são avaliados ao supervisionar as estratégias ; decisões sobre principais transações, processos de gestão de riscos, e politicas correlatas.
     Como se desenvolve o processo de definição de metas e métricas
     

Resumo do IFRS S1 REQUISITOS GERAIS PARA DIVULGAÇÃO DE INFORMAÇÕES FINANCEIRAS RELACIONADAS À SUSTENTABILIDADE
Autoria IFRC
Divulgação junho 2023
Visto ser da missão da IFRS estabelaçer padrões, o documento é normativo e define a forma e o processo para a elaboraçao de relatórios de riscos (positivos e negativos) a sustentabilidade de uma entidade.
Por ser documento normativo tem uma leitura difícil, em especial para leitores não versados em  gestão de riscos e governança corporativa.Ao menos versados no assunto recomento iniciar pela definição de termos do apendice A. 
No que se refere a identificação dos riscos, que o documento referência como ´tópicos de divulgação', o documento faz referência a outras normas:  IFRR S2 - que trata riscos relacionados ao Clima e a normas SASB que orienta a identificação de riscos por setor, alem destas o documento faz referências as especificidades de cada entidade que devem orientar a idnetificação dos ´tópicos de divulgaçaão´
Observa-se no documento os seguintes aspectos:
Definiçao em relaçáo a materialidade , que , ao meu entendimento,trata de como proceder na definiçao de metricas e metas, refere a forma de mensuraçao da relevancia dos tópicos de divulgaçaoa a da forma como mensurar os aspectos econômicos e não econômicos da gestão dos riscos.
Conteúdo Principal
GOVERNANÇA
O documento apresenta requisitos para a apresentaçao da estrutura de governança da sustentabilide, tanto como tema dentro da governança da entidade e como tema próprio que requer por sí uma estrutura de governaná específica.
ESTRATÉGIA
O documento exige a apresentaçao da estratégia em termos de : critérios para a identificaçao dos riscos;  critérios para avaliaçao de impactos; a apresentaçao da Cadeia De Valor da entidade e dentro dessa cadeia onde se concentram os riscos; e  os critérios de domadas de descisões .
Para a apresntaçao do riscos e oportunidades, o documeto define como requisitos a descriçao dos ´tópicos de divulgação' e seu impacto na perpctiva de tempo, sendo assim o documento orienta para a definiáo de curto , médio e longo prazos. Alem de exisgir a mensuraçao dos impactos financeiros no fluxo de caixa da entidade nestas perspecitivas de tempo. 
Alem do aspecto de gerenciamento dos riscos envolvendo a sustenatabilidade 
Assim como a Governaça o documento exige a apresentaçao da ESTRTÉGIA 


